% ROZDZIAŁ 20 — miary Radona i ich regularność
\chapter{Miary borelowskie, regularność i miary Radona}
\label{ch:radon-measures-20}

\providecommand{\atlasref}[2]{#2}


\section{Dlaczego miara potrzebuje topologii?}

\noindent\emph{Powiązania tematyczne w Bibliotece matematycznej:} \atlasref{GT-19::def:topog-polskie-borel}{GT-19: zbiory borelowskie}.

Do tej pory wszystkie miary były określone na przestrzeniach
mierzalnych. Jeśli jednak przestrzeń ma topologię, można pytać,
czy wartości miary zbioru dają się odtworzyć z miar zbiorów
otwartych, domkniętych lub zwartych. Są to trzy różne własności:
regularność zewnętrzna, regularność wewnętrzna i ciasność.
Rozróżnianie ich będzie niezbędne w twierdzeniu reprezentacyjnym
Riesza--Markowa--Kakutaniego.

\begin{definicja}[Miara borelowska i lokalna skończoność]
\label{def:borel-locfinite-20}
Niech $(X,\tau)$ będzie przestrzenią topologiczną,
a $\mathcal B(X)$ jej $\sigma$-algebrą borelowską.
\emph{Miarą borelowską} na $X$ nazywamy nieujemną miarę
$\mu:\mathcal B(X)\to[0,\infty]$.
Miara $\mu$ jest \emph{lokalnie skończona}, jeśli dla
każdego $x\in X$ istnieje otwarte otoczenie $U_x$ punktu $x$
takie, że $\mu(U_x)<\infty$.
\end{definicja}

\begin{uwaga}[Skończoność na zbiorach zwartych]
\label{rem:compact-finiteness-20}
Na dowolnej przestrzeni topologicznej lokalna skończoność miary borelowskiej
implikuje, że $\mu(K)<\infty$ dla każdego zwartego
\emph{borelowskiego} zbioru $K\subseteq X$ (w przestrzeni Hausdorffa każdy zbiór zwarty jest domknięty, a więc borelowski): ze skończonego podpokrycia
$K\subseteq U_{x_1}\cup\cdots\cup U_{x_m}$ dostajemy
$\mu(K)\le\sum_i\mu(U_{x_i})<\infty$.
Na przestrzeni lokalnie zwartej Hausdorffa zachodzi
również implikacja odwrotna: każdy punkt posiada otwarte
otoczenie o zwartym domknięciu, więc jego miara nie
przekracza skończonej miary tego domknięcia.
W przestrzeni, która nie jest przestrzenią Hausdorffa, zbiór zwarty nie musi być borelowski, dlatego samo wyrażenie $\mu(K)$ może nie być określone. W dalszej części rozdziału przestrzenie są Hausdorffa. Bez lokalnej zwartości powyższy dowód implikacji odwrotnej nie ma zastosowania.
\end{uwaga}

\begin{definicja}[Regularność względem topologii]
\label{def:regularity-20}
Niech $X$ będzie przestrzenią Hausdorffa, a $\mu$ miarą
borelowską na $X$. Mówimy, że jest
\emph{zewnętrznie regularna}, gdy dla każdego
$E\in\mathcal B(X)$ zachodzi
\begin{equation}\label{eq:outer-regular-20}
 \mu(E)=\inf\{\mu(U):E\subseteq U,\ U\subseteq X
 \text{ otwarty}\}.
\end{equation}
Mówimy, że jest \emph{wewnętrznie regularna względem
zbiorów domkniętych}, gdy dla każdego $E\in\mathcal B(X)$
\[
 \mu(E)=\sup\{\mu(F):F\subseteq E,\ F\subseteq X
 \text{ domknięty}\}.
\]
Jest \emph{wewnętrznie regularna względem zbiorów zwartych},
gdy dla każdego $E\in\mathcal B(X)$
\begin{equation}\label{eq:inner-compact-20}
 \mu(E)=\sup\{\mu(K):K\subseteq E,\ K\subseteq X
 \text{ zwarty}\}.
\end{equation}
Przyjmujemy $\inf\varnothing=+\infty$ i
$\sup\varnothing=0$ dla podanych rodzin miar nieujemnych.
\end{definicja}

\begin{definicja}[Miara Radona --- przyjęta konwencja]
\label{def:radon-measure-20}
Niech $X$ będzie lokalnie zwartą przestrzenią Hausdorffa.
\emph{Miarą Radona} na $X$ nazywamy miarę borelowską $\mu$,
która jest skończona na zbiorach zwartych, regularna
zewnętrznie na wszystkich zbiorach borelowskich oraz
\emph{regularna wewnętrznie na zbiorach otwartych}, to znaczy
\[
 \mu(U)=\sup\{\mu(K): K\subseteq U,\ K\text{ zwarty}\}
 \qquad(U\subseteq X\text{ otwarty}).
\]
Skończoność na zbiorach zwartych jest na $X$ równoważna
lokalnej skończoności (uwaga~\ref{rem:compact-finiteness-20}).
\end{definicja}

\begin{uwaga}[Dwie konwencje regularności]
W definicji~\ref{def:radon-measure-20} regularność
wewnętrzną wymagamy na zbiorach otwartych, a nie na wszystkich
borelowskich. To rozróżnienie jest konieczne na ogólnych
przestrzeniach lokalnie zwartych Hausdorffa. Silniejsza
regularność na wszystkich zbiorach borelowskich zachodzi,
jeśli przestrzeń jest $\sigma$-zwarta lub jeśli rozważany
zbiór ma miarę $\sigma$-skończoną.
W rozdziale~\ref{ch:haar-23} podamy jawny przykład
miary Haara, dla której silniejszy warunek zawodzi.
Na lokalnie zwartych przestrzeniach metrycznych ośrodkowych
obie konwencje są równoważne.
\end{uwaga}

\begin{lemat}[Regularność na zbiorach $\sigma$-skończonych]
\label{lem:radon-sigmafinite-inner-20}
Niech $X$ będzie lokalnie zwartą przestrzenią Hausdorffa,
a $\mu$ miarą Radona według definicji
\ref{def:radon-measure-20}. Jeżeli $E\in\mathcal B(X)$
i $\mu(E)<\infty$, to
\[
 \mu(E)=\sup\{\mu(K):K\subseteq E,\ K\text{ zwarty}\}.
\]
Równość ta zachodzi także dla każdego borelowskiego
$E$ o mierze $\sigma$-skończonej, czyli dla którego
istnieją borelowskie $E_j\subseteq E$ ze
$\mu(E_j)<\infty$ i $E=\bigcup_jE_j$.
\end{lemat}
\begin{proof}
Niech $\mu(E)<\infty$ i $\varepsilon>0$.
Z regularności zewnętrznej wybieramy otwarty
$U\supseteq E$ z $\mu(U)<\mu(E)+\varepsilon$.
Mamy $\mu(U\setminus E)<\varepsilon$.
Następnie wybieramy otwarty $W\supseteq U\setminus E$
z $\mu(W)<\mu(U\setminus E)+\varepsilon<2\varepsilon$,
a z regularności wewnętrznej na otwartym $U$ zbiór zwarty
$F\subseteq U$ z $\mu(F)>\mu(U)-\varepsilon$.
Wówczas $K=F\setminus W$ jest zwarty jako domknięty
podzbiór zwartego $F$. Ponieważ $F\subseteq U$
oraz $U\setminus E\subseteq W$, mamy $K\subseteq E$.
Ze skończoności $\mu(F)$ i $\mu(W)$ oraz z rozkładu
$F=K\sqcup(F\cap W)$ dostajemy
\[
 \mu(K)=\mu(F)-\mu(F\cap W)
 \ge\mu(F)-\mu(W)>\mu(E)-3\varepsilon.
\]
Dowolność $\varepsilon$ daje nierówność od dołu
dla supremum miar zbiorów zwartych zawartych w $E$,
a nierówność odwrotna wynika z monotoniczności miary.
Dla zbioru o mierze $\sigma$-skończonej zastępujemy dane pokrycie zbiorami skończonej miary jego rosnącymi skończonymi sumami $E_j\subset E$. Jeśli $\mu(E)=\infty$, dla dowolnego $M>0$ wybieramy $j$ z $\mu(E_j)>M+1$ na mocy ciągłości miary z dołu, a następnie zwarty $K\subseteq E_j$ o mierze większej od $\mu(E_j)-1$. Dostajemy $\mu(K)>M$. Jeśli $\mu(E)<\infty$, stosujemy pierwszy przypadek bezpośrednio. W szczególności gdy $X$ jest $\sigma$-zwarta, każdy borelowski $E$ ma miarę $\sigma$-skończoną: przecinamy go z przeliczalnym pokryciem zwartymi zbiorami skończonej miary.
\end{proof}

\section{Regularność skończonych miar borelowskich na przestrzeniach metrycznych}
\begin{twierdzenie}[Aproksymacja zbiorów borelowskich]
\label{thm:finite-metric-regular-20}
Niech $(X,d)$ będzie dowolną przestrzenią metryczną,
a $\mu$ skończoną miarą borelowską na $X$.
Dla każdego $E\in\mathcal B(X)$ i $\varepsilon>0$
istnieją zbiór domknięty $F$ i zbiór otwarty $U$ takie,
że $F\subseteq E\subseteq U$ oraz
\begin{equation}\label{eq:finite-regular-error-20}
 \mu(U\setminus F)<\varepsilon.
\end{equation}
W szczególności $\mu$ jest zewnętrznie regularna
i wewnętrznie regularna względem zbiorów domkniętych.
Nie twierdzimy w tym miejscu, że $F$ można wybrać zwarty.
\end{twierdzenie}
\begin{proof}
\emph{Krok 1: zbiory otwarte.}
Dla zbioru otwartego $U\neq X$ definiujemy
$F_n=\{x\in X:d(x,X\setminus U)\ge1/n\}$.
Funkcja $x\mapsto d(x,X\setminus U)$ jest
$1$-lipschitzowska, więc $F_n$ jest domknięty.
Z otwartości $U$ wynika, że każdy $x\in U$
ma dodatnią odległość od $X\setminus U$;
stąd $F_n\uparrow U$. Ponieważ $\mu(U)<\infty$,
ciągłość miary z dołu daje
$\mu(U\setminus F_n)\to0$. Jeżeli $U=X$,
wystarczy przyjąć $F_n=X$. Zatem każdy zbiór otwarty
spełnia warunek z tezy dla dowolnego $\varepsilon>0$.

\emph{Krok 2: rodzina zbiorów aproksymowalnych.}
Oznaczmy przez $\mathcal D$ rodzinę tych zbiorów
$E\in\mathcal B(X)$, dla których warunek tezy zachodzi
dla każdego $\varepsilon>0$. Rodzina ta zawiera wszystkie
zbiory otwarte. Wykażemy, że jest $\sigma$-algebrą.
Jeżeli $F\subseteq E\subseteq U$ i
$\mu(U\setminus F)<\varepsilon$, to
$X\setminus U\subseteq X\setminus E\subseteq X\setminus F$,
a zbiór $X\setminus U$ jest domknięty,
$X\setminus F$ zaś otwarty. Różnica między nimi
jest równa $U\setminus F$, więc $X\setminus E\in\mathcal D$.

Niech teraz $E=\bigcup_{n\ge1}E_n$, gdzie $E_n\in\mathcal D$.
Dla danego $\varepsilon>0$ dobieramy
$F_n\subseteq E_n\subseteq U_n$, z $F_n$ domkniętym,
$U_n$ otwartym oraz
$\mu(U_n\setminus F_n)<\varepsilon2^{-n-3}$.
Ponieważ $\mu(E)<\infty$ i
$\bigcup_{n\le N}E_n\uparrow E$, istnieje $N$ takie,
że $\mu(E\setminus\bigcup_{n\le N}E_n)<\varepsilon/4$.
Stawiamy $F=\bigcup_{n\le N}F_n$ i $U=\bigcup_{n\ge1}U_n$.
Zbiór $F$ jest domknięty jako suma skończona,
$U$ jest otwarty, a $F\subseteq E\subseteq U$.
Ponieważ
\[
U\setminus F\subseteq
(U\setminus E)\cup(E\setminus F),
\quad
U\setminus E\subseteq\bigcup_{n\ge1}(U_n\setminus E_n),
\]
oraz
$E\setminus F\subseteq
(E\setminus\bigcup_{n\le N}E_n)
\cup\bigcup_{n\le N}(E_n\setminus F_n)$,
podaddytywność daje jawne oszacowanie
\[
 \mu(U\setminus F)\le
 \sum_{n\ge1}\varepsilon2^{-n-3}+
 \varepsilon/4+\sum_{n\le N}\varepsilon2^{-n-3}
 \le\varepsilon/2<\varepsilon.
\]
W konsekwencji $\mathcal D$ jest $\sigma$-algebrą.

\emph{Krok 3: zakończenie.}
Ponieważ $\mathcal D$ zawiera wszystkie zbiory otwarte,
zawiera także $\sigma$-algebrę przez nie generowaną,
czyli $\mathcal B(X)$. Równości w definicjach regularności
wynikają z monotoniczności miary oraz
$\mu(E\setminus F),\mu(U\setminus E)\le\mu(U\setminus F)$.
\end{proof}

\begin{uwaga}[Dlaczego nie można automatycznie zastąpić $F$ zbiorem zwartym?]
W ogólnej przestrzeni metrycznej zbiór domknięty
i ograniczony nie musi być zwarty. Na przykład zamknięta
kula jednostkowa w nieskończenie wymiarowej przestrzeni
Hilberta nie jest zwarta. Istotnie, wybieramy indukcyjnie
ciąg ortonormalny $(e_n)$: na każdym kroku dopełnienie
ortogonalne przestrzeni rozpiętej przez poprzednie wektory
zawiera niezerowy wektor, który normalizujemy. Dla $n\ne m$
mamy $\|e_n-e_m\|^2=2$, więc żaden podciąg tego ciągu
nie jest ciągiem Cauchy'ego, a tym bardziej nie jest zbieżny.
Nawet lokalna zwartość nie wystarcza do zwartości
zbiorów domkniętych i ograniczonych: nieskończony zbiór
z metryką dyskretną $d(x,y)=1$ dla $x\ne y$
jest lokalnie zwarty (singletony są otwarte i zwarte),
domknięty i ograniczony, lecz pokrycie singletonami
nie ma skończonego podpokrycia.
Przykłady wyjaśniają jedynie, dlaczego nie wolno
zastąpić domkniętości i ograniczoności zwartością;
nie wykluczają regularności względem zbiorów zwartych
miar na tych przestrzeniach. Aproksymacja przez
zbiory zwarte wymaga osobnego argumentu.
\end{uwaga}

\section{Ciasność i miary na przestrzeniach polskich}

\noindent\emph{Powiązania tematyczne w Bibliotece matematycznej:} \atlasref{GT-19::def:topog-polskie}{GT-19: przestrzenie polskie}.

\begin{definicja}[Ciasność pojedynczej miary]
\label{def:tight-single-20}
Niech $X$ będzie przestrzenią Hausdorffa, a $\mu$
skończoną miarą borelowską. Mówimy, że $\mu$ jest
\emph{ciasna}, jeżeli dla każdego $\varepsilon>0$
istnieje zwarty zbiór $K\subseteq X$ taki,
że $\mu(X\setminus K)<\varepsilon$.
\end{definicja}

\begin{twierdzenie}[Ciasność skończonych miar na przestrzeniach polskich]
\label{thm:polish-tight-20}
Niech $X$ będzie przestrzenią polską, tj. przestrzenią
topologiczną z metryką zupełną i ośrodkiem przeliczalnym,
zgodną z jej topologią. Każda skończona miara borelowska
$\mu$ na $X$ jest ciasna. Ponadto dla każdego
$E\in\mathcal B(X)$ zachodzi
\[
 \mu(E)=\sup\{\mu(K):K\subseteq E,\ K\text{ zwarty}\}.
\]
\end{twierdzenie}
\begin{proof}
Jeżeli $X=\varnothing$, wybieramy $K=\varnothing$ i teza jest
spełniona. W przeciwnym razie wybieramy zgodną metrykę
zupełną $d$ oraz niepusty co najwyżej przeliczalny
zbiór gęsty $D=\{x_j:j\ge1\}$ (w przypadku skończonym
możemy powtarzać jego elementy). Ustalmy $\varepsilon>0$.
Dla każdego $n\ge1$ kule otwarte
$B(x_j,2^{-n})$, $j\ge1$, pokrywają $X$.
Z ciągłości miary z dołu (twierdzenie \ref{thm:measure-continuity}),
stosowanej do ich skończonych sum, oraz ze skończoności
$\mu(X)$ wynika, że istnieje $m_n$ takie, że
\[
\mu\left(X\setminus\bigcup_{j=1}^{m_n}B(x_j,2^{-n})\right)
<\varepsilon 2^{-n-1}.
\]
Niech $C_n=\bigcup_{j=1}^{m_n}\overline{B}(x_j,2^{-n})$.
Zbiór $C_n$ jest domknięty i
$\mu(X\setminus C_n)<\varepsilon2^{-n-1}$.
Kładziemy $K=\bigcap_{n\ge1}C_n$.
Zbiór $K$ jest domknięty, a dla każdego $n$
zawiera się w sumie skończenie wielu kul o promieniu
$2^{-n}$. Jest zatem całkowicie ograniczony:
dla każdego $\delta>0$ wybieramy $n$ tak duże,
by $2^{1-n}<\delta$, i otrzymujemy skończone
$\delta$-pokrycie $K$ przez kule o środkach w $K$:
z każdej kuli domkniętej w $C_n$, która przecina $K$,
wybieramy punkt $y\in K$; jej przecięcie z $K$ zawiera się
w $B(y,\delta)$ z nierówności trójkąta i $2^{1-n}<\delta$.
Ponieważ $K$ jest domknięty w zupełnym $X$, jest zupełny.
Zupełność i całkowita ograniczoność implikują zwartość:
dla dowolnego ciągu w $K$ wybieramy kolejno nieskończone
podciągi zawarte w jednej z kul skończonego pokrycia
promieniami $2^{-m}$, $m=1,2,\ldots$. Wybór przekątniowy
daje podciąg Cauchy'ego, ponieważ odległość dwóch wyrazów
od $m$-tego miejsca nie przekracza $2^{1-m}$.
Jest on zbieżny w $K$ dzięki zupełności. W przestrzeni
metrycznej zwartość ciągowa jest równoważna zwartości. Wreszcie
\[
 \mu(X\setminus K)
 \le\sum_{n\ge1}\mu(X\setminus C_n)<\varepsilon.
\]
Udowodniliśmy ciasność.

Teraz ustalmy $E\in\mathcal B(X)$. Z twierdzenia
\ref{thm:finite-metric-regular-20} wybieramy domknięty
$F\subseteq E$ taki, że $\mu(E\setminus F)<\varepsilon/2$.
Z ciasności wybieramy zwarty $K_0\subseteq X$
z $\mu(X\setminus K_0)<\varepsilon/2$.
Zbiór $K=F\cap K_0$ jest zwarty, zawiera się w $E$, a
$\mu(E\setminus K)\le\mu(E\setminus F)+\mu(X\setminus K_0)
<\varepsilon$. Monotoniczność daje żądaną równość.
\end{proof}

\begin{przyklad}[Rozkłady gaussowskie w nieskończonym wymiarze]
Twierdzenie~\ref{thm:polish-tight-20} ma zastosowanie
do każdej skończonej miary borelowskiej na ośrodkowej
przestrzeni Banacha z metryką normową, również w nieskończonym
wymiarze. Nie oznacza to, że domknięte kule są zwarte;
zwarty zbiór zapewniający ciasność może mieć znacznie
bardziej skomplikowaną geometrię. W szczególności
przestrzeń $\ell^2$ jest polska, więc jej skończone
miary borelowskie są ciasne.
\end{przyklad}

\section{Miary Radona na przestrzeniach lokalnie zwartych}

\noindent\emph{Powiązania tematyczne w Bibliotece matematycznej:} \atlasref{GT-14::def:topog-baire-lokalna-zwartosc}{GT-14: lokalna zwartość}.

\begin{twierdzenie}[Lokalnie skończone miary borelowskie]
\label{thm:lch-secondcountable-radon-20}
Niech $X$ będzie lokalnie zwartą, ośrodkową
przestrzenią metryczną, a $\mu$ lokalnie skończoną
miarą borelowską na $X$. Wówczas $\mu$ jest miarą
Radona w sensie definicji~\ref{def:radon-measure-20}.
Ponadto $\mu$ jest $\sigma$-skończona.
\end{twierdzenie}
\begin{proof}
\emph{Krok 1: otwarte pokrycie skończonej miary.}
Jeśli $X=\varnothing$, teza zachodzi dla miary zerowej.
W przypadku niepustym skończone podpokrycie można
uzupełnić zbiorami pustymi do ciągu.
Dla każdego $x\in X$ istnieje otwarte otoczenie $W_x$
z $\mu(W_x)<\infty$. Wybieramy otwarte otoczenie $O_x$ o zwartym domknięciu, a następnie promień $r>0$ taki, że $B(x,2r)\subset O_x\cap W_x$. Dla $V_x=B(x,r)$ domknięcie $\overline{V_x}$ zawiera się w $B(x,2r)\subset O_x\cap W_x$ i jako domknięty podzbiór zwartego $\overline{O_x}$ jest zwarte.
W szczególności $\mu(\overline{V_x})<\infty$.
Ponieważ przestrzeń metryczna ośrodkowa spełnia drugi
aksjomat przeliczalności i ma własność Lindelöfa
(por. \atlasref{GT-09::thm:topog-metryczne-trzy-warunki}{GT-09: równoważność ośrodkowości, drugiego aksjomatu przeliczalności i własności Lindelöfa}),
z pokrycia $(V_x)_{x\in X}$ wybieramy przeliczalne
podpokrycie $(V_n)_{n\ge1}$. Dla każdego $n$
$\overline{V_n}$ jest zwarty i ma skończoną miarę.
Już stąd wynika $\sigma$-skończoność $\mu$.
Ponadto każdy zwarty $K\subseteq X$ ma skończoną
miarę: z pokrycia $(V_n)$ wybieramy skończone
podpokrycie $K\subseteq V_{n_1}\cup\dots\cup V_{n_m}$,
a z podaddytywności
\[
 \mu(K)\le\sum_{j=1}^m\mu(V_{n_j})<\infty.
\]
Zwartość w przestrzeni Hausdorffa zapewnia przy tym,
że $K$ jest domknięty, a więc borelowski.

\emph{Krok 2: regularność zewnętrzna.}
Rozbijamy $X$ na parami rozłączne zbiory borelowskie
$B_1=V_1$, $B_n=V_n\setminus\bigcup_{j<n}V_j$.
Dla danego $E\in\mathcal B(X)$ i $\varepsilon>0$
rozważamy skończoną miarę
$\mu_n(A)=\mu(A\cap V_n)$ na przestrzeni metrycznej $X$.
Z twierdzenia~\ref{thm:finite-metric-regular-20}
wybieramy otwarty $U_n\supseteq E\cap B_n$ taki, że
$\mu_n(U_n\setminus(E\cap B_n))<\varepsilon2^{-n-1}$.
Niech $U=\bigcup_{n\ge1}(U_n\cap V_n)$.
Wówczas $U$ jest otwarty, $E\subseteq U$, ponieważ
$E\cap B_n\subseteq U_n\cap V_n$. Zbiory $B_n$ przypisują każdemu punktowi dokładnie jeden indeks pokrycia. Rozłączność nie jest jednak konieczna dla oszacowania regularności: analogicznie można aproksymować $E\cap V_n$, gdyż błędy sumuje się za pomocą podaddytywności.
Ponadto $(U_n\cap V_n)\setminus E\subseteq (U_n\cap V_n)\setminus(E\cap B_n)$, a miara tego drugiego zbioru wynosi $\mu_n(U_n\setminus(E\cap B_n))$. Stąd
\[
\mu(U\setminus E)
\le\sum_{n\ge1}\mu((U_n\cap V_n)\setminus(E\cap B_n))
<\varepsilon.
\]
Jeżeli $\mu(E)<\infty$, daje to
$\mu(U)<\mu(E)+\varepsilon$; jeżeli $\mu(E)=\infty$,
regularność zewnętrzna jest natychmiastowa,
ponieważ każdy otwarty nadzbiór ma miarę nieskończoną.

\emph{Krok 3: regularność wewnętrzna względem zwartych.}
Niech $E\in\mathcal B(X)$ i $\varepsilon>0$.
Dla każdego $n$ rozważamy przestrzeń zwartą
$\overline{V_n}$ z miarą skończoną
$\mu|_{\overline{V_n}}$ i metryką indukowaną.
Zbiór $E\cap V_n$ jest borelowski w podprzestrzeni $\overline{V_n}$, ponieważ $\mathcal B(\overline{V_n})$ jest śladem $\mathcal B(X)$ (otwarte zbiory podprzestrzeni są przecięciami ze zbiorami otwartymi w $X$, a ślad zachowuje dopełnienia względne i przeliczalne sumy). Z twierdzenia~\ref{thm:finite-metric-regular-20}
dla zbioru $E\cap V_n$ wybieramy zbiór
$K_n$ domknięty w $\overline{V_n}$, zawarty w
$E\cap V_n$, o własności
$\mu((E\cap V_n)\setminus K_n)<\varepsilon2^{-n-1}$.
Każdy $K_n$ jest zwarty w $X$.
Dla $E$ o skończonej mierze wybieramy $N$ takie, że
$\mu(E\setminus\bigcup_{n\le N}V_n)<\varepsilon/2$.
Zbiór $K=\bigcup_{n\le N}K_n$ jest zwarty i zawiera się
w $E$, a
\[
 \mu(E\setminus K)
 \le\mu(E\setminus\bigcup_{n\le N}V_n)
 +\sum_{n\le N}\mu((E\cap V_n)\setminus K_n)
 <\varepsilon.
\]
Gdy $\mu(E)=\infty$, z ciągłości miary z dołu
$\mu(E\cap\bigcup_{n\le N}V_n)\uparrow\infty$.
Dla dowolnego $M>0$ wybieramy $N$ z miarą tej części
większą niż $M+1$ i ponawiamy aproksymację przez
$K_n$ dla $1\le n\le N$, z błędami sumującymi się do mniej niż $1$. Zbiór $E\cap\bigcup_{n\le N}V_n$ ma skończoną miarę, a różnica między nim i $K=\bigcup_{n\le N}K_n$ ma miarę mniejszą niż $1$. Dostajemy zwarty $K\subseteq E$ z $\mu(K)>M$.
To dowodzi \eqref{eq:inner-compact-20} także w przypadku
nieskończonej masy. Razem z krokami 1--2 daje tezę.
\end{proof}

\begin{uwaga}[Na jakich przestrzeniach twierdzenie obowiązuje?]
Założenie ośrodkowości i metryczności nie jest częścią
definicji miary Radona; zostało wprowadzone w powyższym
twierdzeniu, aby można było podać samodzielny dowód.
Na ogólnych lokalnie zwartych przestrzeniach Hausdorffa
regularność nie każdej miary borelowskiej jest automatyczna.
Nie wystarczy również samo założenie lokalnej skończoności.
\end{uwaga}

\begin{przyklad}[Lokalna skończoność nie wystarcza]
\label{ex:nonradon-discrete-20}
Niech $X$ będzie nieprzeliczalnym zbiorem z topologią
dyskretną. Jest to lokalnie zwarta przestrzeń metryczna,
ale nieośrodkowa. Wszystkie podzbiory są borelowskie.
Definiujemy
\[
\mu(A)=\begin{cases}0,&A\text{ jest co najwyżej przeliczalny},\\
+\infty,&A\text{ jest nieprzeliczalny}.
\end{cases}
\]
To rzeczywiście miara: $\mu(\varnothing)=0$, a dla
przeliczalnej rodziny zbiorów rozłącznych ich suma jest
nieprzeliczalna wtedy i tylko wtedy, gdy któryś ze zbiorów
jest nieprzeliczalny. W takim przypadku obie strony
równości addytywności są nieskończone, a w przeciwnym
przypadku obie są zerowe.
Miara jest lokalnie skończona, ponieważ każde otwarte
jednoelementowe otoczenie ma miarę zero. Tymczasem
zbiory zwarte w przestrzeni dyskretnej są skończone,
więc wszystkie mają miarę zero, choć $\mu(X)=\infty$.
Regularność wewnętrzna względem zwartych zawodzi.
\end{przyklad}

\section{Przykłady miar Radona}

\begin{przyklad}[Dwie konwencje regularności na nieprzeliczalnej sumie przestrzeni]
\label{it:ex:radon-two-regularities}
Niech $J$ będzie zbiorem nieprzeliczalnym, a
$X=J\times[0,1]$, gdzie $J$ wyposażamy w topologię
dyskretną, a $[0,1]$ w zwykłą topologię.
Przestrzeń $X$ jest lokalnie zwarta i metryzowalna
(na różnych składowych przyjmujemy odległość $2$,
a na tej samej zwykłą odległość w $[0,1]$).
Każda składowa $X_j=\{j\}\times[0,1]$ jest otwarta
i domknięta oraz zwarta, co daje lokalną zwartość.
Metryka opisana powyżej spełnia nierówność trójkąta:
jeśli końce leżą w różnych składowych, co najmniej
jeden odcinek pośredni ma długość $2$; jeśli leżą
w jednej składowej, stosujemy zwykłą nierówność
trójkąta na $[0,1]$ albo zauważamy, że przejście
przez inną składową ma długość $4$.
Kule o promieniu mniejszym od $1$ generują topologię
na każdej składowej, więc metryka daje zadaną topologię.
Przestrzeń nie jest $\sigma$-zwarta: każdy zwarty podzbiór
przecina tylko skończenie wiele składowych, ponieważ
składowe tworzą otwarte pokrycie $X$. Przeliczalna
suma zwartych zbiorów przecina więc co najwyżej
przeliczalnie wiele składowych.

Dla $E\in\mathcal B(X)$ oznaczmy
$E_j=\{t\in[0,1]:(j,t)\in E\}$ oraz
$J(E)=\{j\in J:E_j\ne\varnothing\}$.
Każdy $E_j$ jest borelowski jako przeciwobraz
$E$ przez ciągłe odwzorowanie $t\mapsto(j,t)$.
Określmy dwie funkcje zbioru:
\[
 \nu(E)=\sum_{j\in J}\lambda_1(E_j),\qquad
 \mu(E)=
 \begin{cases}
 \nu(E),& J(E)\text{ jest przeliczalny},\\
 +\infty,& J(E)\text{ jest nieprzeliczalny}.
 \end{cases}
\]
Sumę nieprzeliczalnej rodziny liczb nieujemnych
rozumiemy jako supremum sum po skończonych
podrodzinach.

Funkcja $\nu$ jest miarą. Istotnie, dla rozłącznych
$E_n$ i skończonego $F\subseteq J$ mamy
\begin{align*}
 \sum_{j\in F}\lambda_1((\bigcup_n E_n)_j)
 &=\sum_n\sum_{j\in F}\lambda_1((E_n)_j)\\
 &\le\sum_n\nu(E_n).
\end{align*}
Supremum po $F$ daje jedną nierówność.
Dla każdego $N$ i skończonych $F_1,\ldots,F_N\subseteq J$
mamy
$\sum_{n=1}^N\sum_{j\in F_n}\lambda_1((E_n)_j)
\le\nu(\bigcup_n E_n)$, stosując skończoną sumę
po $F_1\cup\cdots\cup F_N$ i addytywność na składowych.
Niezależne supremum po $F_n$ i następnie po $N$
daje nierówność odwrotną.
Funkcja $\mu$ również jest miarą: jeżeli choć jeden
$J(E_n)$ jest nieprzeliczalny, obie strony warunku
addytywności są nieskończone. W przeciwnym razie
$J(\bigcup_n E_n)=\bigcup_nJ(E_n)$ jest przeliczalny,
więc addytywność wynika z addytywności $\nu$.
Obie miary są lokalnie skończone, ponieważ każda
składowa $X_j$ jest otwarta i ma masę $1$.

Dla otwartego $U\subseteq X$ zachodzi $\mu(U)=\nu(U)$.
Gdy $J(U)$ jest nieprzeliczalny, wszystkie niepuste
$U_j$ są względnie otwarte i mają dodatnią miarę
Lebesgue'a. Ponieważ rodzina dodatnich liczb o skończonej
sumie ma co najwyżej przeliczalnie wiele wyrazów
(dla każdego $n\ge1$ tylko skończenie wiele może być
większych od $1/n$), w tym przypadku $\nu(U)=\infty$.
Dla przeliczalnego $J(U)$ równość wynika z definicji.

Miara $\mu$ jest zewnętrznie regularna. Jeśli
$\mu(E)=\infty$, wynika to z monotoniczności.
Jeżeli $\mu(E)<\infty$, to $J(E)$ jest przeliczalny;
numerujemy jego elementy i dla każdego z nich
(z numeracją rozpoczynającą się od $1$) aproksymujemy $E_j$ od góry zbiorem względnie otwartym
$U_j$ w $[0,1]$ o błędzie mniejszym niż
$\varepsilon2^{-n}$ dla $j$ o numerze $n$,
korzystając z twierdzenia
\ref{thm:finite-metric-regular-20}.
Suma otwartych zbiorów $\{j\}\times U_j$ zawiera $E$
i ma miarę $\mu$ mniejszą od $\mu(E)+\varepsilon$
(dla $E=\varnothing$ wybieramy zbiór pusty).
Miara $\mu$ jest też wewnętrznie regularna na otwartych:
dla skończonego $F\subseteq J$ przybliżamy każde $U_j$,
$j\in F$, zwartym $K_j\subseteq U_j$ na mocy
twierdzenia \ref{thm:finite-metric-regular-20}
na zwartym $[0,1]$, z błędem mniejszym niż
$\varepsilon/(1+|F|)$ na każdej składowej. Skończona suma
$K=\bigcup_{j\in F}(\{j\}\times K_j)$ jest zwarta
i $\mu(K)=\sum_{j\in F}\lambda_1(K_j)$.
Dla każdego skończonego $F$ mamy więc
$\mu(K)>\sum_{j\in F}\lambda_1(U_j)-\varepsilon$.
Dowolność błędu oraz supremum po $F$ dają
$\sup_{K\subseteq U,\ K\text{ zwarty}}\mu(K)
=\nu(U)=\mu(U)$, także dla nieskończonej masy.
Skończoność na zwartych wynika z lokalnej skończoności.
Zatem $\mu$ jest miarą Radona według definicji
\ref{def:radon-measure-20}.

Rozważmy jednak domknięty zbiór $E=J\times\{0\}$.
Mamy $\mu(E)=+\infty$, natomiast każdy zwarty
$K\subseteq E$ jest skończony i $\mu(K)=0$.
Regularność wewnętrzna względem zwartych nie zachodzi
więc na wszystkich zbiorach borelowskich.
Z kolei $\nu(E)=0$, lecz każdy otwarty nadzbiór $U$
zbioru $E$ ma $\nu(U)=+\infty$, zatem $\nu$ nie jest
zewnętrznie regularna, mimo lokalnej skończoności.
Obie miary zgadzają się na całkach funkcji z $C_c(X)$,
ponieważ ich nośniki przecinają tylko skończenie
wiele składowych, na których miary są identyczne.
To pokazuje znaczenie regularności w twierdzeniu
o jednoznaczności miary Radona.
\end{przyklad}

\begin{przyklad}[Miara Lebesgue'a jako miara Radona]
Na $\mathbb R^d$, $d\ge1$, ograniczenie miary Lebesgue'a
$\lambda_d$ do $\mathcal B(\mathbb R^d)$ jest miarą
lokalnie skończoną: każdy punkt ma otoczenie będące
kulą o skończonej objętości. Przestrzeń $\mathbb R^d$
jest lokalnie zwarta, metryczna i ośrodkowa.
Na mocy twierdzenia~\ref{thm:lch-secondcountable-radon-20}
miara ta jest Radona. Jej uzupełnienie na
$\mathcal L_d$ zawiera także zbiory nieborelowskie;
trzeba rozróżniać tę miarę zupełną od jej ograniczenia
borelowskiego.
\end{przyklad}

\begin{przyklad}[Miary z lokalnie całkowalną gęstością]
Niech $w:\mathbb R^d\to[0,\infty]$ będzie borelowska
i lokalnie całkowalna względem $\lambda_d$,
tj. $\int_Kw\,\mathrm d\lambda_d<\infty$ dla każdego
zwartego $K\subseteq\mathbb R^d$. Wówczas
$\mu_w(E)=\int_Ew(x)\,\mathrm d\lambda_d(x)$
dla $E\in\mathcal B(\mathbb R^d)$ określa miarę na mocy
lematu \ref{lem:integration-density-18}. Warunek lokalnej
całkowalności gwarantuje lokalną skończoność, ponieważ
kula otwarta o skończonym promieniu zawiera się w zwartej
kuli domkniętej. Regularność wynika z twierdzenia
\ref{thm:lch-secondcountable-radon-20}. Na przykład $w(x)=\exp(|x|^2)$ daje miarę
Radona o nieskończonej masie całkowitej.
\end{przyklad}

\begin{przyklad}[Atomy gęste i nośnik całej prostej]
Niech $(q_n)_{n\ge1}$ będzie numeracją zbioru
$\mathbb Q$ bez powtórzeń. Na $\mathbb R$ określmy
$\mu=\sum_{n\ge1}2^{-n}\delta_{q_n}$.
Jest to skończona miara borelowska, a więc miara Radona.
Każdy niepusty przedział otwarty zawiera pewien $q_n$
i ma dodatnią miarę $\mu$, choć $\mu(\mathbb R\setminus
\mathbb Q)=0$. W konsekwencji nośnikiem topologicznym
$\mu$ jest cała prosta. Pełny nośnik topologiczny
nie oznacza bezwzględnej ciągłości względem miary Lebesgue'a:
w tym przykładzie $\mu\perp\lambda_1$.
\end{przyklad}

\section{Nośnik topologiczny}
\begin{definicja}[Nośnik miary]
\label{def:support-radon-20}
Dla miary borelowskiej $\mu$ na przestrzeni topologicznej
$X$ określamy jej \emph{nośnik topologiczny} przez
\[
\operatorname{supp}\mu
:=\{x\in X:\mu(U)>0\text{ dla każdego otwartego }
U\ni x\}.
\]
Jest to zbiór domknięty. Nie jest to z definicji
wybrany mierzalny zbiór skupiający miarę w sensie
rozdziału~\ref{ch:ac-singular-17}.
\end{definicja}

\begin{twierdzenie}[Nośnik na przestrzeni spełniającej drugi aksjomat przeliczalności]
\label{thm:support-fullmass-20}
Niech $X$ będzie przestrzenią spełniającą drugi aksjomat przeliczalności,
a $\mu$ dowolną miarą borelowską na $X$. Wtedy
$\operatorname{supp}\mu$ jest domknięty,
$\mu(X\setminus\operatorname{supp}\mu)=0$ i jest
najmniejszym domkniętym zbiorem o pełnej mierze
w sensie zawierania.
\end{twierdzenie}
\begin{proof}
Z definicji
$X\setminus\operatorname{supp}\mu$
jest sumą wszystkich otwartych zbiorów o mierze zero,
a więc jest otwarty. Niech $(B_n)$ będzie przeliczalną
bazą topologii. Jeżeli $x$ należy do tej sumy, to
istnieje otwarty $U\ni x$ o $\mu(U)=0$.
Można wybrać $B_n$ takie, że $x\in B_n\subseteq U$.
Wówczas $\mu(B_n)=0$, a więc
\[
X\setminus\operatorname{supp}\mu
=\bigcup_{\{n:\mu(B_n)=0\}}B_n.
\]
Przeliczalna podaddytywność daje miarę zero tej sumy.
Niech $F\subseteq X$ będzie domknięty i
$\mu(X\setminus F)=0$. Każdy $x\notin F$ ma
otwarte otoczenie $X\setminus F$ miary zero,
zatem $x\notin\operatorname{supp}\mu$.
W konsekwencji $\operatorname{supp}\mu\subseteq F$.
\end{proof}

\begin{uwaga}[Dlaczego założenie przeliczalności ma znaczenie?]
W przykładzie~\ref{ex:nonradon-discrete-20} każda
jednopunktowa okolica ma miarę zero, zatem
$\operatorname{supp}\mu=\varnothing$, choć
$\mu(X)=+\infty$. Problem wynika z faktu, że
nieprzeliczalna suma zbiorów zerowych nie musi być
zerowa. Pokazuje to, dlaczego nie należy dowodzić
pełności miary nośnika przez nieuzasadnione przejście
od dowolnego pokrycia do przeliczalnej sumy.
Dla miary Radona w przyjętej konwencji drugi aksjomat
przeliczalności nie jest jednak potrzebny: otwarty zbiór
$U=X\setminus\operatorname{supp}\mu$ jest sumą otwartych
zbiorów zerowych. Każdy zwarty $K\subseteq U$ ma skończone
podpokrycie takimi zbiorami, zatem $\mu(K)=0$.
Regularność wewnętrzna na otwartych daje $\mu(U)=0$.
Argument minimalności z poprzedniego dowodu pozostaje ważny.
\end{uwaga}

\section{Funkcje ciągłe o zwartym nośniku}
Dla lokalnie zwartej przestrzeni Hausdorffa $X$ oznaczamy
przez $C_c(X)$ przestrzeń ciągłych funkcji
$f:X\to\mathbb R$, których \emph{nośnik funkcji}
$\operatorname{supp}f=\overline{\{x:f(x)\neq0\}}$
jest zwarty. W analogiczny sposób definiuje się
$C_c(X;\mathbb C)$. Nie należy mylić nośnika funkcji
z nośnikiem topologicznym miary.

\begin{lemat}[Funkcja odcinająca]
\label{lem:cutoff-radon-20}
Niech $X$ będzie lokalnie zwartą przestrzenią Hausdorffa,
$K\subseteq X$ zwartym zbiorem, a $U\subseteq X$
otwartym zbiorem zawierającym $K$. Wtedy istnieje
$f\in C_c(X)$ taka, że $0\le f\le1$, $f|_K=1$
oraz $\operatorname{supp}f\subseteq U$.
\end{lemat}
\begin{proof}
Jeśli $K=\varnothing$, wystarczy przyjąć $f=0$.
W dalszej części zakładamy $K\neq\varnothing$.
Z lematu \atlasref{GT-14::lem:topog-baire-zwarte-domkniecie}{GT-14: otoczenie o zwartym domknięciu} dla każdego $x\in K$
wybieramy otwarte $W_x\ni x$ o zwartym domknięciu zawartym
w $U$. Ze zwartości $K$ otrzymujemy skończone podpokrycie
$W_{x_1},\ldots,W_{x_m}$. Dla
$V=\bigcup_{i=1}^m W_{x_i}$ mamy
$K\subseteq V\subseteq\overline V=\bigcup_{i=1}^m\overline{W_{x_i}}\subseteq U$,
a $\overline V$ jest zwarty jako skończona suma zwartych zbiorów.

Przestrzeń $\overline V$ jest normalna na mocy twierdzenia
\atlasref{GT-07::thm:topog-zwartosc-hausdorff-normalna}{GT-07: normalność zwartych przestrzeni Hausdorffa}.
Zbiory $K$ i $\overline V\setminus V$ są w niej domknięte i rozłączne.
Z lematu \atlasref{GT-05::thm:topog-urysohn}{GT-05: lemat Urysohna} dla przestrzeni normalnych otrzymujemy ciągłą funkcję $g:\overline V\to[0,1]$, równą $1$ na $K$ oraz $0$ na $\overline V\setminus V$ (gdy drugi zbiór jest pusty, można przyjąć $g=1$). Przedłużamy $g$ zerem na $X\setminus V$. Definicje zgadzają się na przecięciu dwóch zbiorów domkniętych $\overline V$ i $X\setminus V$, które pokrywają $X$, więc funkcja $f$ jest ciągła. Jej nośnik zawiera się w $\overline V\subseteq U$, co daje tezę.

W przypadku metrycznym można zamiast lematu Urysohna użyć jawnego wzoru: jeżeli $V=X$, przyjmujemy $f=1$, a w przeciwnym razie określamy
\[
 f(x)=\frac{d(x,X\setminus V)}
 {d(x,X\setminus V)+d(x,K)}.
\]
Mianownik nie zeruje się nigdzie: dla $x\in K$
pierwszy składnik jest dodatni, a poza $K$ drugi
składnik jest dodatni (ponieważ $K$ jest domknięty).
Obie funkcje odległości są ciągłe, stąd $f$ jest
ciągła; na $K$ ma wartość $1$, poza $V$ jest zerowa,
a $\operatorname{supp}f\subseteq\overline V\subseteq U$.
\end{proof}

\begin{wniosek}[Charakteryzacja nośnika przez funkcje testowe]
\label{cor:support-test-20}
Niech $X$ będzie lokalnie zwartą przestrzenią Hausdorffa,
a $\mu$ miarą Radona na $X$. Dla każdego $x\in X$
następujące warunki są równoważne:
\begin{enumerate}[label=\textup{(\roman*)}]
\item $x\in\operatorname{supp}\mu$;
\item dla każdej funkcji $f\in C_c(X)$ takiej, że
$f\ge0$ i $f(x)>0$, zachodzi
$\int_X f(y)\,\mathrm d\mu(y)>0$.
\end{enumerate}
\end{wniosek}
\begin{proof}
Załóżmy (i) i wybierzmy $f$ jak w (ii).
Z ciągłości $f$ wynika, że istnieją $\delta>0$
oraz otwarte otoczenie $V$ punktu $x$ takie,
że $f(y)\ge\delta$ dla $y\in V$. Zatem
$\int_Xf\,\mathrm d\mu\ge\delta\mu(V)>0$.
Całka jest skończona, ponieważ $f$ jest
ograniczona i ma zwarty nośnik, a $\mu$
jest skończona na zbiorach zwartych.

Jeśli natomiast $x\notin\operatorname{supp}\mu$,
to istnieje otwarty zbiór $U\ni x$ z $\mu(U)=0$.
Lemat~\ref{lem:cutoff-radon-20}, zastosowany do
$K=\{x\}\subseteq U$, daje $f\in C_c(X)$
spełniającą $f\ge0$, $f(x)=1$ i
$\operatorname{supp}f\subseteq U$.
Stąd $\int_X f\,\mathrm d\mu=0$, co przeczy (ii).
\end{proof}

\begin{twierdzenie}[Jednoznaczność miary Radona przez całki]
\label{thm:uniqueness-cc-20}\label{tw:radon-unikalnosc}
Niech $X$ będzie lokalnie zwartą przestrzenią Hausdorffa,
a $\mu,\nu$ miarami Radona na $X$. Jeżeli
\[
\int_X f(x)\,\mathrm d\mu(x)
=\int_Xf(x)\,\mathrm d\nu(x)
\qquad\text{dla każdego }f\in C_c(X),
\]
to $\mu=\nu$ na $\mathcal B(X)$.
\end{twierdzenie}
\begin{proof}
Dla każdego otwartego $U\subseteq X$ i dowolnej
miary Radona $\rho$ na $X$ zachodzi
\begin{equation}\label{eq:cc-open-sup-20}
\rho(U)=\sup\left\{\int_Xf\,\mathrm d\rho:
 f\in C_c(X),\ 0\le f\le1,\
 \operatorname{supp}f\subseteq U\right\}.
\end{equation}
Istotnie, każda z powyższych całek jest co najwyżej
$\rho(U)$. Z drugiej strony dla każdego zwartego
$K\subseteq U$ lemat~\ref{lem:cutoff-radon-20}
daje $f$ równe $1$ na $K$, więc
$\int_Xf\,\mathrm d\rho\ge\rho(K)$.
Z regularności wewnętrznej miary $\rho$
uzyskujemy \eqref{eq:cc-open-sup-20} także dla
$\rho(U)=+\infty$. Przy założonej równości całek
zbiory funkcji w supremum są te same dla obu miar,
a więc $\mu(U)=\nu(U)$ dla wszystkich otwartych $U$.
Wreszcie dla dowolnego zbioru borelowskiego
$E\subseteq X$ z regularności zewnętrznej
$\mu(E)=\inf_{U\supseteq E}\mu(U)
=\inf_{U\supseteq E}\nu(U)=\nu(E)$.
\end{proof}

\begin{uwaga}[Związek z twierdzeniem reprezentacyjnym]
Twierdzenie~\ref{thm:uniqueness-cc-20} daje
\emph{jednoznaczność} miary reprezentującej funkcjonał
$L(f)=\int_Xf\,\mathrm d\mu$ na $C_c(X)$.
Istnienie takiej miary dla dodatniego funkcjonału
jest przedmiotem następnego rozdziału, poświęconego
twierdzeniu Riesza--Markowa--Kakutaniego. W przestrzeni
niezwartej funkcjonał całkowania względem miary
Radona o nieskończonej masie nie musi być ograniczony
w normie supremum na całym $C_c(X)$; dlatego rozróżnimy
wtedy $C_c(X)$ i $C_0(X)$.
\end{uwaga}

\begin{twierdzenie}[Gęstość $C_c(X)$ w przestrzeniach $L^p$]
\label{thm:cc-dense-lp-radon-20}
Niech $X$ będzie lokalnie zwartą przestrzenią Hausdorffa,
a $\mu$ miarą Radona w sensie definicji
\ref{def:radon-measure-20}. Wówczas dla każdego
$1\le p<\infty$ przestrzeń $C_c(X;\mathbb C)$ jest
gęsta w $L^p(X,\mathcal B(X),\mu;\mathbb C)$.
Analogiczny wniosek zachodzi dla funkcji rzeczywistych.
\end{twierdzenie}
\begin{proof}
Najpierw ustalmy $E\in\mathcal B(X)$ takie, że
$\mu(E)<\infty$. Dla $\varepsilon>0$ z lematu \ref{lem:radon-sigmafinite-inner-20}
(stosowanego do zbioru skończonej miary $E$) i regularności
zewnętrznej wybieramy zwarty $K$
i otwarty $U$ takie, że
$K\subseteq E\subseteq U$ oraz
$\mu(E\setminus K)<\varepsilon/2$,
$\mu(U\setminus E)<\varepsilon/2$.
Lemat~\ref{lem:cutoff-radon-20} daje $g\in C_c(X)$
taką, że $0\le g\le1$, $g|_K=1$ oraz
$\operatorname{supp}g\subseteq U$.
Na $K$ funkcje $\chi_E$ i $g$ są równe, poza $U$
obie znikają, a w pozostałej części ich różnica
ma moduł co najwyżej $1$. Zatem
\[
 \int_X|\chi_E(x)-g(x)|^p\,\mathrm d\mu(x)
 \le\mu(E\setminus K)+\mu(U\setminus E)
 <\varepsilon.
\]
Otrzymujemy aproksymację funkcji charakterystycznej
każdego zbioru o skończonej mierze.

Z twierdzenia \ref{thm:simple-dense-lp-14} (które nie
wymaga $\sigma$-skończoności miary) dla każdego $f\in L^p(\mu)$
i $\eta>0$ wybieramy funkcję prostą
$s=\sum_{j=1}^m a_j\chi_{E_j}$, gdzie $\mu(E_j)<\infty$,
z $\|f-s\|_p<\eta/2$. Jeżeli wszystkie $a_j$ są zerowe,
wystarczy funkcja zerowa. W przeciwnym razie z poprzedniego
akapitu wybieramy $g_j\in C_c(X)$ tak, aby
$\|\chi_{E_j}-g_j\|_p<\eta/(2(1+\sum_{k=1}^m|a_k|))$.
Nierówność Minkowskiego (twierdzenie \ref{thm:minkowski-14}) daje
\[
 \left\|f-\sum_{j=1}^m a_jg_j\right\|_p
 \le \|f-s\|_p+\sum_{j=1}^m|a_j|\|\chi_{E_j}-g_j\|_p
 <\eta.
\]
Suma $\sum_{j=1}^m a_jg_j$ jest więc żądaną aproksymacją. Jest ona ciągła, a jej nośnik zawiera się w skończonej sumie
zwartych nośników funkcji $g_j$, więc jest zwarty.
To dowodzi żądanej gęstości.
\end{proof}

\begin{uwaga}[Znaczenie założenia $p<\infty$]
Gęstość w normie $L^p$ nie rozciąga się bez dodatkowych
warunków na $p=\infty$. Na przykład względem miary
Lebesgue'a na $[0,1]$ funkcji $\chi_{[0,1/2]}$
nie da się aproksymować w normie $L^\infty$
funkcjami ciągłymi. Istotnie, gdyby dla ciągłej $g$
zachodziło $\|g-\chi_{[0,1/2]}\|_\infty<1/2$, wybralibyśmy
$r<1/2$ większe od tej normy. Wówczas $|g(x)-1|\le r$
na całym $(0,1/2)$ oraz $|g(x)|\le r$ na całym $(1/2,1)$:
naruszenie którejkolwiek nierówności w punkcie, przez ciągłość,
zachodziłoby na otwartym przedziale dodatniej miary.
Przechodząc do granicy w $1/2$ z obu stron, otrzymujemy
$1\le|1-g(1/2)|+|g(1/2)|\le2r<1$, sprzeczność.
\end{uwaga}

\begin{lemat}[Regularność miar zdominowanych]
\label{it:lem:radon-dominated}
Niech $X$ będzie lokalnie zwartą przestrzenią Hausdorffa.
Jeżeli $\sigma$ jest skończoną miarą Radona na $X$,
a $\tau$ nieujemną miarą borelowską z $\tau\le\sigma$,
to $\tau$ jest skończoną miarą Radona.
Skończona suma skończonych miar Radona jest również
skończoną miarą Radona.
\end{lemat}
\begin{proof}
Dla borelowskiego $E$ i $\varepsilon>0$ regularność
$\sigma$ i lemat \ref{lem:radon-sigmafinite-inner-20}
dają zwarty $K\subseteq E$ oraz otwarty $U\supseteq E$
z $\sigma(U\setminus K)<\varepsilon$ (wybieramy
błędy aproksymacji od wewnątrz i z zewnątrz
mniejsze niż $\varepsilon/2$). Wtedy
$\tau(U\setminus K)<\varepsilon$. Ponieważ
$\tau(X)\le\sigma(X)<\infty$, możemy odejmować
skończone wartości i otrzymujemy
$\tau(K)>\tau(E)-\varepsilon$ oraz
$\tau(U)<\tau(E)+\varepsilon$. Dowolność
$\varepsilon$ daje obie regularności dla $\tau$. Dla skończonej sumy $m$ miar wybieramy
dla każdej z nich analogiczną parę $K_i,U_i$ z
błędem mniejszym niż $\varepsilon/m$. Kładąc
$K=\bigcup_iK_i$, $U=\bigcap_iU_i$, dostajemy
zwarty podzbiór $E$, otwarty nadzbiór $E$ oraz
$U\setminus K\subseteq U_i\setminus K_i$ dla każdego $i$.
Zsumowanie błędów daje tezę.
\end{proof}

\section{Regularność a uzupełnienie miary}
Jeżeli $\mu$ jest miarą borelowską, jej uzupełnienie
(wprowadzone w rozdziale~2) może zawierać zbiory
nieborelowskie, ale nie zmienia wartości miary
na zbiorach borelowskich. Na przykład miara Lebesgue'a
na $\mathcal L_d$ jest uzupełnieniem miary borelowskiej
$\lambda_d|_{\mathcal B(\mathbb R^d)}$. Z drugiej strony
nie każda miara Radona jest zupełna na swojej
$\sigma$-algebrze borelowskiej: jeśli borelowski zbiór
zerowy ma podzbiór nieborelowski, to uzupełnienie
ma większą dziedzinę. Jest to rozróżnienie
istotne przy formułowaniu twierdzeń o obrazach miar,
całkowaniu i odsyłaczach do $L^p$.

\section{Co przygotowaliśmy na następny rozdział?}
Przeprowadziliśmy dwa kierunki regularności:
miary skończone na przestrzeniach metrycznych są
aproksymowane od wewnątrz przez zbiory domknięte,
a na przestrzeniach polskich przez zbiory zwarte;
miary lokalnie skończone na lokalnie zwartych
przestrzeniach metrycznych ośrodkowych są miarami Radona.
Ustaliliśmy ponadto, że wszystkie całki z funkcji
$C_c(X)$ jednoznacznie wyznaczają miarę Radona.
Następny rozdział odwróci tę konstrukcję:
zacznie od dodatniego funkcjonału na funkcjach ciągłych
i skonstruuje reprezentującą go miarę.
